\chapter{Project Defense and Project Expo}
\label{chap:project-defense-expo}

\noindent
This chapter presents the project evaluation and dissemination processes conducted through
two rounds of project defense and participation in the SKE Project Expo. This chapter
documents the feedback received from evaluators, the revisions made in response to that
feedback, and the outcomes of those revisions. In addition, it discusses the team's experience
in presenting the project to a broader audience, including poster design, audience interaction,
and reflections on how the Expo contributed to the overall development of the project.

\section{First Project Defense}
\label{sec:first-defense}

\subsection{Comments and Feedback from the Evaluators}
\label{subsec:first-defense-feedback}

During the first project defense, the evaluators observed that the project was primarily presented from a research-oriented perspective. While this direction was valid at the early stage of development, it led to some ambiguity regarding the overall purpose and positioning of the project.

In particular, the evaluators found that the intended target users were not clearly defined. It was unclear whether the system was designed for researchers, developers, or general users. This made it difficult to determine how the system would be used in practice and what its main contribution was.

The evaluators also suggested that the team should provide a clearer explanation of what the project is doing and what research problem it is addressing. Although technical details were presented, especially in relation to EEG analysis and machine learning, the connection between these components and the overall objective was not sufficiently clear. As a result, both the research motivation and the system’s role were not communicated effectively.

Furthermore, the system interface appeared highly technical and resembled a “cockpit-style” design, which reduced its accessibility. This reinforced the perception that the project was more research-focused rather than a user-oriented software system.

Overall, the main concern raised during the first defense was that the project was still too research-oriented in its presentation, with unclear target users and insufficient explanation of its practical purpose.

\subsection{Revisions Required, Actions Taken, and Outcomes}
\label{subsec:first-defense-revisions}

Following the first project defense, the team recognized that the main issue was the lack of clarity in project positioning, particularly in terms of target users and overall purpose. However, at this stage, the project was still in an early phase of development, and its research-oriented nature was still appropriate.

Instead of making immediate structural changes, the team focused on reflecting on the feedback and reassessing the direction of the project. This included considering how the system could evolve from a research prototype into a more complete and user-oriented software platform, as well as how to better communicate its objectives and contributions.

As a result, the first defense served as an important turning point by highlighting the gap between technical implementation and clear communication. Although no major revisions were implemented immediately, the insights gained during this stage directly informed the significant changes made prior to the second project defense.

\section{Second Project Defense}
\label{sec:second-defense}

\subsection{Comments and Feedback from the Evaluators}
\label{subsec:second-defense-feedback}

In the second project defense, the evaluators observed that the project had improved in terms of overall clarity and structure. The revised presentation made it easier to understand the purpose of the system and how its components are connected. The project was also recognized as being conceptually valid and implementable, indicating that the core idea had been successfully demonstrated.

However, the evaluators noted that the presentation was still relatively technical and research-heavy, which could limit accessibility for a broader audience. While the direction of the project was clearer than in the first defense, further improvements in simplifying explanations and emphasizing practical usage were recommended.

In addition, the evaluators raised questions regarding the system architecture, data processing pipeline, and potential real-world applications. These questions reflected a growing interest in understanding not only the technical components, but also how the system operates as a complete solution and how it could be applied in practical contexts.

\subsection{Additional Revisions, Actions Taken, and Final Outcomes}
\label{subsec:second-defense-revisions}

The main revision involved redefining the project as an interactive EEG analysis platform, rather than solely a machine learning experiment. This shift led to several concrete changes.

First, the system design was adjusted to emphasize modularity and usability, allowing users to interact more effectively with preprocessing, visualization, and modeling components. Second, the presentation narrative was restructured to highlight the overall workflow and system architecture, instead of focusing primarily on model-level details.

In addition, the target users of the system were more clearly defined as researchers and practitioners working with EEG data, which helped align both the system design and the presentation with practical use cases.

As a result of these changes, the project became more aligned with software engineering objectives. The system's purpose---supporting reproducible EEG analysis and experimentation---was communicated more clearly. This revision significantly improved the coherence of the project and demonstrated a more mature and well-defined direction.

\subsection{Suggestions on How Project Defense Can Better Support Software Development}
\label{subsec:defense-suggestions}

The project defense process played an important role in guiding the development of the system, particularly in clarifying the project's direction and improving its communication. Feedback from the evaluators helped the team recognize that technical implementation alone is not sufficient; clear explanation of objectives, target users, and system usage is equally important.

However, the process could be further improved to better support software development. For example, providing earlier feedback on project scope and positioning would help reduce ambiguity during the initial stages. Additionally, clearer evaluation criteria that distinguish between research contributions and software engineering aspects would allow students to better align their work with expectations.

Furthermore, emphasizing usability, system architecture, and real-world applicability during evaluation would encourage teams to develop more practical and user-oriented solutions. Strengthening these aspects would make the defense process more effective in supporting the development of complete and impactful software systems.

\section{Project Expo Participation}
\label{sec:expo-participation}

\subsection{Poster Presentation and Poster Components}
\label{subsec:expo-poster}

The poster presented at the Project Expo was designed to communicate the project as an integrated EEG analysis platform. It was structured into several key sections, including project overview, objectives, system architecture, data pipeline, platform visualization, study setup, predictive models, research findings, and future work.

The main objective of the poster was to convey, within a short time, that the project provides a unified and reproducible workflow for EEG data analysis and prestimulus behavioral prediction. To achieve this, the design emphasized a balance between technical content and visual clarity. Diagrams illustrating the system architecture and data pipeline were used to provide an intuitive understanding of how the platform operates, while concise text descriptions summarized the key components and findings.

The poster also highlighted both aspects of the project: the software platform and the underlying machine learning models. This ensured that the audience could understand not only how the system functions, but also what insights it produces.

\begin{figure}[H]
    \centering
    \includegraphics[width=0.9\textwidth]{figures/eeg_poster.jpg}
    \caption{Poster}
    \label{fig:eeg-poster}
\end{figure}

\subsection{Event Photos and Awards Received}
\label{subsec:expo-photos}

During the Project Expo, the team actively participated by presenting the system at their assigned booth. A live demonstration of the platform was conducted, allowing visitors to observe the workflow and interact with the system in real time. This approach helped communicate the practical functionality of the project more effectively.

The team did not receive any awards during the event. However, the experience of presenting and demonstrating the system provided valuable opportunities for engagement and feedback from a diverse audience.

\begin{figure}[H]
    \centering
    \includegraphics[width=0.9\textwidth]{figures/expo_pic.jpg}
    \caption{Event Photo}
    \label{fig:expo-photo}
\end{figure}

\subsection{Questions Raised During the Expo and the Team's Responses}
\label{subsec:expo-questions}

During the Expo, visitors raised several questions regarding the technical details and practical applications of the project. One common question concerned the source of the data used in the system. The team explained that the platform utilizes EEG data from the Healthy Brain Network dataset, which is widely used in research.

Another frequently asked question was about the type of artificial intelligence used. The team described the use of EEGNet, a convolutional neural network designed specifically for EEG signal analysis, and explained its role in both regression and classification tasks.

Visitors also asked about the system architecture and data processing pipeline. In response, the team explained the workflow, including data preprocessing, model execution, and visualization, using diagrams provided in the poster.

In addition, questions were raised about potential real-world applications of the system. The team discussed possible use cases such as behavioral prediction, cognitive state analysis, and support tools for EEG-based research.

These interactions highlighted the importance of clearly explaining both the technical and practical aspects of the system, and they provided valuable external perspectives on the project.

\subsection{Recommendations for Improving the Expo}
\label{subsec:expo-recommendations}

Overall, the Project Expo was well-organized and provided a valuable platform for presenting student projects. The event facilitated interaction between teams and a broad audience, enabling knowledge sharing and constructive discussions.

While no major issues were identified, there are opportunities for improvement. Increasing the involvement of industry representatives could provide additional insights into real-world applications and enhance the practical relevance of the projects. Additionally, more structured feedback mechanisms could help participants better understand their strengths and areas for improvement.